% Related Works is Appendix A, and the same file carries every other appendix
% section, each opened by \clearpage with an apdx: label [venue: ICML]. The
% file is named after its content (latex_source.md section 11).
% Floats here use [ht]/[!ht].
\section{Related Works} \label{apdx:related-works}
% 3 subsections x one paragraph (100-170 words each).
\subsection{<Base Paradigm>}
<Field> is still hard for current models \cite{Placeholder-A}. Existing approaches mostly belong to <A> \cite{Placeholder-B} or <B> \cite{Placeholder-C}. Both, however, share <shared limitation>.

\subsection{<Enabling Technique>}
<Technique> has started to attract attention \cite{Placeholder-A, Placeholder-B}. Although useful, these methods are prone to <limitation 1> and <limitation 2>.

\subsection{<Form of Supervision>}
On the one hand, <line 1> \cite{Placeholder-B}. On the other hand, <line 2> \cite{Placeholder-C}. This progress motivates our design of <component>.

\clearpage
\section{Implementation Details} \label{apdx:impl-details}
All experiments use the <framework>\footnote{<framework URL>} framework.
Table~\ref{tab:hyper} lists the key hyperparameters.

\begin{table}[ht]
\centering
\caption{Key hyperparameters.}
\label{tab:hyper}
\begin{tabular}{l|c}
\toprule
\textbf{Hyperparameter} & \textbf{Value} \\
\midrule
Learning rate & \pend \\
Batch size & \pend \\
Random seeds & \pend \\
\bottomrule
\end{tabular}
\end{table}

\clearpage
\section{Prompts} \label{apdx:prompts}
Figure~\ref{fig:full-prompt} shows the full instruction given to <module> of X.
\begin{figure}[ht]
    \centering
    \begin{gptpromptbox}
    \raggedright\small\ttfamily
    <full prompt text>
    \end{gptpromptbox}
    \caption{Full prompt of X.}
    \label{fig:full-prompt}
\end{figure}

\clearpage
\section{Experiments} \label{apdx:experiments}
% Subsections mirror the RQs. Paired panels at 0.49\columnwidth (= 3.3 in here).
\subsection{Results Across <Sweep>} \label{apdx:sweep}
Figure~\ref{fig:apdx-sweep} shows the results across <sweep>.
\begin{figure}[!ht]
    \centering
    \begin{subfigure}[t]{0.49\columnwidth}
        \centering
        \figplaceholder[3cm]{Model-L}
        \caption{Model-L.}
    \end{subfigure}
    \begin{subfigure}[t]{0.49\columnwidth}
        \centering
        \figplaceholder[3cm]{Model-S}
        \caption{Model-S.}
    \end{subfigure}
    \caption{Results across <sweep> on <task>.}
    \label{fig:apdx-sweep}
\end{figure}

\clearpage
\section{Case Study} \label{apdx:case-study}
This section walks through several examples of X in action.
